\section{The cutoff polylogarithm and its disappearing depth}
\label{sec:polylog}

The ordered endpoint evaluation is accompanied by an explicit cutoff
identity. This is useful both as a proof of the endpoint cancellation
and as a finite transform for noninteger orders, one of the questions
raised in \cite{repo:finite}.

\subsection{An identity entire in the two orders}

For $|z|<1$ put
\begin{align}
 G_{s,t}(z)&=\sum_{0\le m<n}(n-m)^2
             w_m(s,1)w_n(t,1)z^{n+1},\label{eq:cutoff-g}\\
 B_{s,t}(z)&=\sum_{k\ge1}\frac{z^k}{k^s(k+1)^t},
 \qquad c_s=\frac{1-2^{-s}}2,\qquad \Theta=z\frac{\dd}{\dd z},
 \label{eq:cutoff-b}\\
 U_t(z)&=\frac{(1-z)^2}{z^2}\Li_t(z)+2-z^{-1}-2^{-t}.
 \label{eq:cutoff-u}
\end{align}
The last expression has a removable singularity at zero and equals
$2\sum_{n\ge0}w_n(t,1)z^{n+1}$. The same convention will apply to
the removable negative powers of $z$ in the next formula.

\begin{theorem}[Finite cutoff formula]\label{thm:cutoff}
For all $s,t\in\C$ and $|z|<1$,
\begin{align}
 G_{s,t}(z)={}&\frac{(1-z)^2}{2z^2}\Li_{t,s}(z,1)
 +\frac{z-1}{2z}\Li_{s+t}(z)
 +\frac14B_{t,s}(z)-\frac1{4z}B_{s,t}(z)\notag\\
 &+\frac12\left(c_s\Theta^2-\Theta-\frac12\right)U_t(z)
 +\frac12-2^{-t-2}.
 \label{eq:cutoff-master}
\end{align}
Every term is entire in $s,t$ on the open unit disk. At positive
integer orders the two $B$ functions are finite combinations of
ordinary polylogarithms. The only depth-two term in this representation
has the explicit factor $(1-z)^2$.
\end{theorem}

\begin{proof}
The inner finite identity of Section~\ref{sec:ordered}, with $a=1$,
becomes
\begin{equation}\label{eq:inner-harmonic}
 \sum_{m=0}^{n-1}(n-m)^2w_m(s,1)
 =H_{n+1}^{(s)}+c_s(n+1)^2-(n+1)-\frac12
   +\frac1{2(n+2)^s}.
\end{equation}
It also holds at $n=0$, since the right side then vanishes.
For $k\ge1$, the three generating identities needed for its harmonic
part are
\begin{align}
 \sum_{k\ge1}\frac{z^kH_k^{(s)}}{k^t}
   &=\Li_{t,s}(z,1)+\Li_{s+t}(z),\notag\\
 \sum_{k\ge1}\frac{z^kH_k^{(s)}}{(k+1)^t}
   &=z^{-1}\Li_{t,s}(z,1),\notag\\
 \sum_{k\ge1}\frac{z^kH_k^{(s)}}{(k+2)^t}
   &=z^{-2}\Li_{t,s}(z,1)-z^{-1}B_{s,t}(z).
 \label{eq:harmonic-generators}
\end{align}
They follow by writing the finite harmonic sum and changing the
outer index; the term subtracted in the last line is precisely the
last inner index in the shifted nested sum.

Multiply \eqref{eq:inner-harmonic} by $w_n(t,1)z^{n+1}$ and sum.
Its quadratic polynomial contributes
$\tfrac12(c_s\Theta^2-\Theta-\tfrac12)U_t$.
The harmonic contribution, by \eqref{eq:harmonic-generators}, is
\[
 \frac{(1-z)^2}{2z^2}\Li_{t,s}(z,1)
 +\frac12\Li_{s+t}(z)-\frac1{2z}B_{s,t}(z).
\]
The final term $1/(2(n+2)^s)$ contributes
\[
 \frac14B_{t,s}(z)-\frac1{2z}\Li_{s+t}(z)
 +\frac1{4z}B_{s,t}(z)+\frac12-2^{-t-2}.
\]
Adding gives \eqref{eq:cutoff-master}. On compact sets of orders,
every summand grows at most polynomially in its indices before the
exponential factor is applied. All rearrangements therefore converge
normally and prove the identity for every pair of complex orders.
\end{proof}

\subsection{Finite classical polylogarithms at integer orders}

For positive integers $p,q$, use the coefficients of
\eqref{eq:pf-a}--\eqref{eq:pf-b} with $D=1$. Then
\begin{equation}\label{eq:b-classical}
 B_{p,q}(z)=\sum_{k=1}^{p}A_k\Li_k(z)
 +\sum_{l=1}^{q}B_lz^{-1}\{\Li_l(z)-z\}.
\end{equation}
This follows by summing the rational partial fractions against $z^n$.
At $|z|<1$ the two simple-pole sums need not be paired for convergence,
although pairing is necessary before taking their individual endpoint
values. Since $\Theta\Li_\nu=\Li_{\nu-1}$, the differential term in
\eqref{eq:cutoff-master} is also finite in classical polylogarithms.
Thus the integer-order generator is explicitly in depth at most two,
and its depth-two coefficient is known exactly.

\begin{corollary}[Endpoint cancellation]\label{cor:cutoff-endpoint}
For positive integers $p\ge1$, $q\ge2$, the radial limit is
\begin{equation}\label{eq:cutoff-endpoint}
 \Ord(p,q;1,1)=c_p\zeta(q)-\frac12+2^{-p-2}
 +\frac14\{B_{q,p}(1)-B_{p,q}(1)\}.
\end{equation}
Together with \eqref{eq:b-classical}, interpreted through its paired
simple poles, this yields the parity formula
\eqref{eq:ordered-parity} at $a=1$.
\end{corollary}

\begin{proof}
Here $\Li_{q,p}(1,1)$ converges: the inner harmonic sum is bounded
for $p>1$ and grows only logarithmically for $p=1$, whereas $q\ge2$.
Its factor $(1-z)^2$ makes its contribution zero. The factor
$z-1$ similarly kills $\Li_{p+q}(z)$.
Finite summation by parts for the three moments of the single weight
gives
\[
 U_q(1)=1-2^{-q},\qquad
 (\Theta U_q)(1)=1,\qquad
 (\Theta^2U_q)(1)=2\zeta(q)-1.
\]
These moments converge for $q>1$. Substitution into
\eqref{eq:cutoff-master} gives \eqref{eq:cutoff-endpoint}.
Absolute convergence of the original double sum justifies its radial
limit as well.
\end{proof}

There is a broader analytic interpretation. If $\Rea t>1$ and
$\Rea(s+t)>0$, the original $G_{s,t}$ has the endpoint
$\Ord(s,t;1,1)$. The depth-two contribution still tends to zero
radially. Indeed
$\sum_{j<k}j^{-\Rea s}$ is
$O(1+k^{\max(1-\Rea s,0)}\log(k+1))$, and a direct comparison of
the resulting power series with a Laplace integral proves
$(1-z)^2\Li_{t,s}(z,1)\to0$ on this domain. Some of the remaining
separated terms can diverge and must be kept in the combination in
\eqref{eq:cutoff-master}. The identity of the collective endpoint
with the boundary formula is valid; replacing a divergent component
by an unannounced regularized value is unnecessary.

\subsection{A finite noninteger-order transform}

Writing $B_{s,t}$ as a one-index series already reduces depth, but a
standard special-function transform makes the answer more explicit.
Let
\[
 \Phi(z,v,c)=\sum_{k\ge0}\frac{z^k}{(k+c)^v},\qquad |z|<1,
\]
be the Lerch transcendent \cite{dlmf:lerch}. For $\Rea s,\Rea t>0$,
\begin{equation}\label{eq:b-beta-lerch}
 B_{s,t}(z)=\frac{z\Gamma(s+t)}{\Gamma(s)\Gamma(t)}
 \int_0^1u^{s-1}(1-u)^{t-1}
       \Phi(z,s+t,2-u)\dd u.
\end{equation}
To prove it, apply the beta identity
\begin{equation}\label{eq:feynman-basic}
 X^{-s}Y^{-t}=\frac{\Gamma(s+t)}{\Gamma(s)\Gamma(t)}
 \int_0^1\frac{u^{s-1}(1-u)^{t-1}}{[uX+(1-u)Y]^{s+t}}\dd u
\end{equation}
with $X=k$, $Y=k+1$ and sum $z^k$. The initial exchange is absolutely
convergent. Both sides then continue in their order variables from
that domain; the continued $B$ itself is entire for $|z|<1$.
Equations \eqref{eq:cutoff-master} and \eqref{eq:b-beta-lerch}
give a finite description of the cutoff generator for noninteger
orders using ordinary polylogarithms, one depth-two polylogarithm,
and two ordinary beta--Lerch transforms. No assumption that arbitrary
complex orders have finite integer partial fractions is made.

\subsection{Euler primitives}

For a function vanishing at zero, define
$\mathcal J f(z)=\int_0^z f(u)\dd u/u$ within the unit disk.
Termwise integration yields, for every integer $\ell\ge0$,
\begin{equation}\label{eq:euler-primitives}
 \mathcal J^\ell G_{s,t}(z)
 =\sum_{0\le m<n}\frac{(n-m)^2}{(n+1)^\ell}
             w_m(s,1)w_n(t,1)z^{n+1}.
\end{equation}
At positive integer orders these remain finite combinations of
depth-at-most-two polylogarithms and rational functions. One constructive
proof uses \eqref{eq:inner-harmonic}: the coefficient is a rational
function of $k=n+1$, times $H_k^{(p)}$, plus a rational function.
Dividing by $k^\ell$ introduces only additional integer-order poles
at $0,-1,-2$. Partial fractions reduce each term to one of
\eqref{eq:harmonic-generators} with shifted orders, or to its finite
index-shift analogues. Shifting $H_{k+c}^{(p)}$ changes it by a
finite rational sum, so no further nested index is introduced.
This proves closure under these primitives constructively; it does
not assert a minimal transcendental basis at the endpoint.
